Fine-tuning moves a point. Weight space $\Theta$ is drawn as a chart around the pretrained point $\theta_0$; the grid is the lattice of a round-to-nearest quantiser $\kappa$ through the zero weight $0$ (marked $+$). Solid curves are image germs of methods pointed at $\theta_0$ (\thmref{Definition}{I.2}), drawn darkest near $\theta_0$ because only the germ there matters. LoRA \citep{hu2021lora} is the cone $\theta_0+\Mr$ with its apex at $\theta_0$ (\thmref{Theorem}{II.2}), shown as the rank-$\le1$ double cone $x^2=y^2+z^2$ of the symmetric $2\times2$ slice in perspective. OFT \citep{qiu2023controlling} is the orbit of $\theta_0$ under rotations, an arc of the circle about $0$ through $\theta_0$, the rest of which is dotted. (IA)$^3$ \citep{liu2022few} is a torus orbit, whose closure is the line through $0$ and $\theta_0$ (\thmref{Theorem}{II.5}). BitFit \citep{benzaken2021bitfit} moves along a coordinate line, parallel to the grid. The dashed curves are not pointed at $\theta_0$. Prefix tuning \citep{li2021prefix} is an extension of the network (\thmref{Definition}{VI.1}); drawn in weight space it is a curve that misses $\theta_0$, at the distance marked by the dotted red segment. QLoRA \citep{dettmers2023qlora} is LoRA pointed at the nearest grid node $\kappa\theta_0$, so its cone is LoRA's cone translated by the defect $e=\kappa\theta_0-\theta_0$ (red arrow), and $\theta_0$ is off its image unless $\rk e\le r$ (\thmref{Proposition}{V.2}).
